\section*{A Splittability and $\tau$-Sampling}
Our decoding result follows from the work of Jernoimo, Srivastava, and Tulsiani[20]. In[20], for the result of Theorem35 to hold, they require$\mathbf W$ to be$\tau$-splittable. $\tau$-splittability is stronger than, and in fact implies,$\tau$-sampling. We will not define splittability here, since for their result to hold, only a ``Splittable Mixing Lemma'' is required. The mixing property used in[20] goes as follows.
\paragraph{Definition56 (splittable mixing for $\pm1$ cut functions).} Given positive integers$n,t$, and$s<t-1$, we denote by$\mathcal F_s$ the set of all functions$f:[n]^t\to\{1,-1\}$ of the form
\[f(x_1,\ldots,x_t)=b\cdot\chi_{A_1}(x_1)\cdots\chi_{A_s}(x_s)\cdot\chi_B(x_{s+1},\ldots,x_t),\tag{27}\]
where$b\in\{1,-1\}$, $A_1,\ldots,A_s\subseteq[n]$, $B\subseteq[n]^{t-s}$, and where$\chi_S(x)=-1$ if$x\in S$ and1 otherwise.
For a distribution$\mathbf W\sim[n]^t$, we denote by$\mathbf W_{[i,j]}$ the distribution over$[n]^{j-i+1}$ that is obtained by sampling$\mathbf x\sim\mathbf W$ and outputting$\mathbf x_{[i,j]}$. For simplicity, we abbreviate$\mathbf W_i=\mathbf W_{[i,i]}$. Given$s\in[t-1]$, let$\nu_s$ be the probability measure over$[n]^t$ for which$\nu_s(x)=\prod_{i\in[s]}\Pr[\mathbf W_i=x_i]\cdot\Pr[\mathbf W_{[s+1,t]}=x_{[s+1,t]}]$. We say$\mathbf W\sim[n]^t$ satisfies splittable mixing with error$\tau$, if for every$s\in[t-1]$ and every$f,f'\in\mathcal F_s$ it holds that
\[\left|\mathbb E_{\mathbf x\sim\nu_s}[f(\mathbf x)f'(\mathbf x)]-\mathbb E_{\mathbf x\sim\nu_{s-1}}[f(\mathbf x)f'(\mathbf x)]\right|\leq\tau.\tag{28}\]
In[20], they show that the above property follows from$\tau$-splittablility. Here we show that the above property also follows from our weaker notion of$\tau$-sampling. For the sake of being compatible with[20], we will use a slightly different definition of$\tau$-sampling than what was given in Definition32.
\paragraph{Definition57.} We say that a distribution$\mathbf W$ over$[n]^t$ is$\tau$-sampling if for any$i\in[t-1]$, $S\subseteq[n]$, and$X\subseteq[n]^{t-i}$,
\[\Cov_{\mathbf w\sim\mathbf W}[1[\mathbf w_i\in S],1[(\mathbf w_{i+1},\ldots,\mathbf w_t)\in X]]\leq\tau.\]
For decoding lifted sum codes, both definitions are essentially the same, since we can always work with$\mathbf W_{\mathrm{rev}}$, wherein we sample$\mathbf w=(\mathbf w_1,\ldots,\mathbf w_t)\sim\mathbf W$ and output$(\mathbf w_t,\ldots,\mathbf w_1)$. Moreover, our$\mathbf W$ satisfies$\mathbf W=\mathbf W_{\mathrm{rev}}$.
\paragraph{Claim58.} Let$\mathbf W\subseteq[n]^t$ be homogeneous and$\tau$-sampling. Then,$\mathbf W$ satisfies splittable mixing with error$4\tau$.
\begin{proof}
Let$s\in[t-1]$, $b,b'\in\{1,-1\}$ and sets$A_1,\ldots,A_s,A'_1,\ldots,A'_s\subseteq[n]$ and$B,B'\subseteq[n]^{t-s}$ that correspond to the functions$f$ and$f'$. For$i\in[s]$ denote$C_i=A_i\triangle A'_i$ and$D=B\triangle B'$. We then have
\[f(x)f'(x)=\sigma\cdot\chi_{C_1}(x_1)\cdots\chi_{C_s}(x_s)\cdot\chi_D(x_{s+1},\ldots,x_t),\]
for$\sigma=bb'\in\{1,-1\}$. The subtraction on the left hand side of Equation(28) now reads
\begin{multline*}
\sum_{x\in[n]^{s-1}}\prod_{i\in[s-1]}\Pr[\mathbf W_i=x_i]\cdot\sigma\cdot\chi_{C_1}(x_1)\cdots\chi_{C_{s-1}}(x_{s-1})\cdot\\
\sum_{y\in[n],z\in[n]^{t-s}}\chi_{C_s}(y)\cdot\chi_D(z)\\
\cdot(\Pr[\mathbf W_s=y]\Pr[\mathbf W_{[s+1,t]}=z]-\Pr[\mathbf W_{[s,t]}=y\circ z]).\tag{29}
\end{multline*}
The second line of the above expression can be written as
\begin{multline*}
\mathbb E_{\mathbf y\sim\mathbf W_s}[(-1)^{1[\mathbf y\in C_s]}]\cdot\mathbb E_{\mathbf z\sim\mathbf W_{[s+1,t]}}[(-1)^{1[\mathbf z\in D]}]\\
-\mathbb E_{(\mathbf y,\mathbf z)\sim(\mathbf W_s,\mathbf W_{[s+1,t]})}[(-1)^{1[\mathbf y\in C_s]}(-1)^{1[\mathbf z\in D]}],\tag{30}
\end{multline*}
which is simply the covariance between the random variables$(-1)^{1[\mathbf y\in C_s]}$ and$(-1)^{1[\mathbf z\in D]}$ where$\mathbf y$ and$\mathbf z$ are drawn appropriately. As$\mathbf W$ is$\tau$-sampling, we know that
\[|\Cov[1[\mathbf y\in C_s],1[\mathbf z\in D]]|\leq\tau,\]
and so Equation(30), in absolute value, amounts to
\[|\Cov[(-1)^{1[\mathbf y\in C_s]},(-1)^{1[\mathbf z\in D]}]|\leq4\tau,\]
which readily follows from the fact that$(-1)^b=1-2b$ for$b\in\{0,1\}$, as we argued in Lemma34. Taking absolute value, by the triangle inequality, Equation(29) can be bounded by
\[4\tau\cdot\sum_{x\in[n]^{s-1}}\prod_{i\in[s-1]}\Pr[\mathbf W_i=x_i].\]
As$\mathbf W$ is homogeneous,$\Pr[\mathbf W_i=x_i]=1/n$, and we are done.
\end{proof}
